%% P17 \dash{} Rachunek macierzowy
%%
%% Zalążek. Poniższy opis jest kontraktem tego programu: co ma uzasadnić i co
%% ma dać czytelnikowi. Napisanie programu oznacza usunięcie bloku
%% \programstub{} i zastąpienie go programem. Jeśli po przerobieniu materiału
%% opis okaże się błędny, popraw go i odnotuj sprzeczność w CLAUDE.md w sekcji
%% *Resolved questions*.

\program{Rachunek macierzowy}
\label{prog:P17}

\programstub{%
\textit{[Opis do przetłumaczenia \dash{} patrz wydanie angielskie.]}

Argument: differentiating with respect to a vector or a matrix is
bookkeeping plus a layout convention, and most of the pain in the literature
is the convention rather than the mathematics. Payoff: the identities the
reader actually needs, derived once each: the gradient of \code{Wx} with
respect to \code{W}; of a squared error; of a softmax; of a log-softmax; of
layer normalisation. The headline: the gradient of cross-entropy through a
softmax is \code{p - y}, the single most reused fact in applied machine
learning, derived in full \dash{} and the reason the two operations are fused in
every serious implementation. Also a versionbox on numerator versus
denominator layout, because the reader will meet both in the same week.

\medskip
\textbf{Planowana długość:} 60 ramek.
}
